\chapter{Аналитический раздел}

% \section{Методы рендеринга трехмерных сцен}
% \section{Анализ предметной области} % Камера, сцена, свет


\section{Трассировка лучей}
Трассировка лучей (от англ. Ray tracing) --- это метод рендеринга, который позволяет реалистично имитировать освещение среды за счет использования физически корректной модели поведения света. Распространяясь по сцене и взаимодействую с ее объектами, свет может отражаться, преломляться, поглощаться и рассеиваться.

\subsection{Алгоритм прямой трассировки лучей}
Главная идея, лежащая в основе алгоритма прямой трассировки лучей, заключается в том, что наблюдатель видит объекты посредством испускаемого источниками света, который многократно отражаясь, согласно законам оптики, доходит до глаз наблюдателя. Метод прямой трассировки предполагает построения траекторий лучей от всех источников освещения ко всем точкам всех объектов сцены, лучи отражаются, преломляются или проходят сквозь него и в результате достигает наблюдателя. Если объект не является отражающим или прозрачным, то траектория луча на этой точке обрывается.
\begin{figure}[H]
	\centering
	\includegraphics[width=0.8\textwidth]{forward_ray_tracing.png}
	\caption{Алгоритм прямой трассировки лучей}
\end{figure}
Основным недостатком этого алгоритма является излишне большое число рассматриваемых лучей, приводящее с существенным затратам вычислительных мощностей, так как лишь малая часть лучей достигает точки наблюдения. Однако данный алгоритм имитирует реальное поведение света и применяется в задачах, в которых точность важнее скорости (оптика и фотоника, расчет освещенности помещения и тд).

\subsection{Алгоритм обратной трассировки лучей}
На отображение сцены влияют только лучи, достигшие наблюдателя. Благодаря тому, что свет имеет особенность распространяться одинаково в обоих направлениях, можно искусственно поменять направления лучей на противоположное, и испускать лучи из точки наблюдения. Это позволяет рассматривать исключительно лучи, которые гарантированно попадут в камеру. Данный метод, основанный на испускании лучей от камеры к объектам сцены, называется обратной трассировкой лучей. Данный подход позволяет существенно сократить количество обрабатываемых лучей и обеспечить более эффективное формирование изображений при приемлемой степени физической достоверности.

\begin{figure}[H]
	\centering
	\includegraphics[width=0.8\textwidth]{backward_ray_tracing.pdf}
	\caption{Алгоритм обратной трассировки лучей}
\end{figure}

\subsection{Выбор алгоритма}

\begin{tabular}{|c|c|c|}
	\hline
	Алгоритм & Быстродействие & Реалистичность \\ 
	\hline
	Прямая трассировка лучей & Низкое & Высокая \\ 
	Обратная трассировка лучей & Среднее & Среднее \\ 
	\hline
\end{tabular}

В качестве алгоритма рендеринга был выбран метод обратной трассировки лучей, так как для поставленной задачи не требуется полная физическая корректность модели, а также важна производительность, учитывая, что программа будет выполняться на центральном, а не графическом процессоре.


\section{Анализ моделей освещения}
Свет при взаимодействии с поверхностью может быть отражен, поглощен или преломлен. Часть световой энергии поглощается поверхностью и преобразуется в тепло. Отраженный и преломленный свет делают объект видимым для наблюдателя. Если весь падающий свет будет поглощен поверхностью объекта, он будет невидимым (абсолютное черное тело).

Все модели освещения делятся на два вида: глобальные и локальные. Глобальные модели учитывают возможности отражения и преломления света от объектов, не являющихся прямыми источниками освещения, поэтому они требуют значительных вычислительных затрат. Например, прозрачные поверхности должны позволять просматривать объекты за ними, на поверхностях возможно появление таких явлений, как рефлексы. Глобальные модели освещения основываются на том, что видимость и затенения связаны между собой: яркость точки поверхности определяется распределением яркости по всем остальным поверхностям, видимым из этой точки. Так можно моделировать шероховатость и свойства отражения реальных материалов, освещение многократно отраженным светом и связанные с ним цветовые эффекты.

Существуют более простые, локальные модели освещения, которые учитывают только свет от источника. Они вычисляют интенсивность, цвет и дальнейшее распределение света на поверхности, но не учитывают перенос света между объектами. Изображение формируется в результате отражения, падающего на поверхность объектов света, интенсивность и цвет которого необходимо рассчитать. Фоновое освещение определяет цвет объекта в тени или при отсутствии явных источников освещения. Его интенсивность постоянна, а также равномерно распределена по всему пространству.
Выделяют две основные модели локального освещения: модель Ламберта и модель Фонга.

\subsection{Модель освещения Ламберта}
Модель Ламберта описывает идеальную "матовую" диффузную поверхность. Интенсивность отраженного света подчиняется закону косинусов Ламберта, согласно которому отраженное излучение одинаково во всех направлениях. При этом положение и угол обзора наблюдателя не имеют значения, т.к. диффузно отраженный свет рассеивается равномерно во всех направлениях.

\begin{figure}[H]
	\centering
	\includegraphics[width=0.8\textwidth]{lambert.pdf}
	\caption{Схема диффузного отражения Ламберта}
\end{figure}


\subsection*{Модель освещения Фонга}
\subsection*{Модель освещения Блинна-Фонга}


\subsection{Взаимодействие света с поверхностью}
\section{Обзор геометрических примитивов}

